% !TeX program = xelatex
\documentclass[UTF8,zihao=-4,fontset=none]{ctexart}
\usepackage[a4paper,left=20mm,right=20mm,top=21mm,bottom=20mm,headheight=12pt,headsep=5mm]{geometry}
\usepackage{fontspec,amsmath,amssymb,array,tabularx,booktabs,graphicx,fancyhdr,caption,pgfplots}
\setmainfont{Times New Roman}
\setCJKmainfont{SimSun}[BoldFont=SimHei]
\setCJKfamilyfont{zhsong}{SimSun}[BoldFont=SimHei]
\setCJKfamilyfont{zhhei}{SimHei}
\providecommand{\songti}{\CJKfamily{zhsong}}
\providecommand{\heiti}{\CJKfamily{zhhei}}
\usepgfplotslibrary{groupplots,statistics}
\pgfplotsset{compat=1.18}
\usepackage[hidelinks]{hyperref}
\hypersetup{pdftitle={3×3 2048自主游戏算法实验报告},pdfauthor={蔡明辰、潘亮、陈达、孟维祎}}
\ctexset{section={format=\heiti\fontsize{14bp}{18bp}\selectfont,beforeskip=10bp,afterskip=6bp},
subsection={format=\heiti\fontsize{12bp}{15bp}\selectfont,beforeskip=7bp,afterskip=5bp}}
\linespread{1}
\newcommand{\bodyfont}{\songti\fontsize{12bp}{12bp}\selectfont}
\AtBeginDocument{\bodyfont}
\setlength{\parindent}{2em}
\setlength{\parskip}{5bp}
\setlength{\tabcolsep}{5bp}
\renewcommand{\arraystretch}{1.12}
\newcolumntype{Y}{>{\raggedright\arraybackslash}X}
\newcolumntype{Z}{>{\centering\arraybackslash}X}
\DeclareCaptionFont{report}{\songti\fontsize{12bp}{14bp}\selectfont}
\captionsetup{font=report,labelsep=quad,skip=5bp,hypcap=false}
\renewcommand{\tablename}{表}
\renewcommand{\figurename}{图}
\pagestyle{fancy}\fancyhf{}
\fancyhead[L]{\songti\fontsize{9bp}{10bp}\selectfont 3×3 2048自主游戏算法实验}
\fancyfoot[C]{\songti\fontsize{9bp}{10bp}\selectfont\thepage}
\renewcommand{\headrulewidth}{0.4pt}
\setlength{\emergencystretch}{2em}
\setcounter{topnumber}{3}
\setcounter{bottomnumber}{2}
\renewcommand{\topfraction}{0.9}
\renewcommand{\bottomfraction}{0.8}
\renewcommand{\textfraction}{0.1}
\setlength{\textfloatsep}{10bp}
\setlength{\floatsep}{8bp}
\raggedbottom
\begin{document}
{\centering\heiti\fontsize{18bp}{24bp}\selectfont 3×3 2048自主游戏算法实验报告\par}
{\centering 蔡明辰\quad 潘亮\quad 陈达\quad 孟维祎\par}

\noindent\textbf{摘要：}针对3×3棋盘空间有限、随机落块容易导致死局的问题，本实验比较MCTS、完整Expectimax、TD和PPO在2048中的表现，并考察三种TD价值表示。以终局数字总和为主要指标，在20个固定测试种子下，完整Expectimax的平均总和为761.0，乐观5-tuple TD为612.1，两者均达到20局生成256的结果。搜索方法能够直接利用环境规则，TD则把经验保存在权重中，降低每步决策所需的计算。不同方法的训练量和搜索开销差异较大，因此本实验主要比较已有实现的效果，并讨论结果受计算预算和超参数影响的原因。

\section{任务与环境}
本实验要求智能体在3×3的2048棋盘上自主决策，使终局所有数字之和尽量大，同时尽可能生成不小于256的方块。动作包括左、右、上、下。每次有效移动后，环境从空格中均匀选取一个位置，以0.9和0.1的概率分别生成2和4；没有合法动作时游戏结束。生成256后继续游戏，以避免提前终止影响总和比较。

合并只改变方块的分布，不改变总和。若以$S_t$表示第$t$步的棋盘总和，基础奖励设为$r_t=S_{t+1}-S_t$，有效移动的奖励即为新落块值2或4。提高终局总和需要让棋盘保持可移动，从而容纳更多新方块。

实现上保留原环境的移动、合并与显示功能，通过继承环境类调整棋盘大小、随机源、奖励和统计信息。算法分别实现动作选择接口，配置读取、权重加载和评测由公共入口完成。这样可以在沿用原有代码结构的同时，复用相同规则比较不同策略。

\section{方法}
\subsection{MCTS与完整Expectimax}
MCTS通过反复模拟估计动作价值。搜索树使用UCT选择动作：
\begin{equation}
Q(s,a)+c\sqrt{\frac{\ln(N(s)+1)}{N(s,a)}},
\end{equation}
其中$N(s)$和$N(s,a)$分别为状态与动作的访问次数，$c$控制探索强度。模拟显式考虑随机落块，达到设定深度后，用棋盘总和、空格数、蛇形排列和达标奖励估值，再向根节点回传。其优点是无需离线训练，但有限模拟次数会带来较大的估计误差。

Expectimax利用已知落块概率求期望。玩家节点选择价值最大的动作，机会节点枚举空格以及出生块2和4。记$x$为移动后、落块前的后状态，$A(s)$为合法动作集合，$E(x)$为空格集合，有
\begin{equation}
\begin{aligned}
V(s)&=\max_{a\in A(s)}C(\operatorname{move}(s,a)),\\
C(x)&=\frac{1}{|E(x)|}\sum_{i\in E(x)}[0.9V(x\oplus2_i)+0.1V(x\oplus4_i)].
\end{aligned}
\end{equation}
本实验使用完整搜索，不设置有限深度的启发式截断，终局效用为
\begin{equation}
U(s)=S(s)+256\,\mathbf{1}\{\max(s)\ge256\}.
\end{equation}
因此，算法优化的是期望总和与达标奖励的加权目标。求解器用C++实现，借助棋盘编码、旋转与镜像对称以及缓存减少重复求解。完整搜索有较高的首次求解和存储成本。

\subsection{TD价值学习}
TD根据相邻状态的价值差更新估计。使用后状态可以把玩家移动与随机落块分开：对每个合法动作，先计算移动后的棋盘，再比较其长期价值。以未来总和增量为价值时，TD(0)误差为
\begin{equation}
\delta=z+\gamma V(x')-V(x),
\end{equation}
其中$z$为新落块值，终局后续价值为0。乐观5-tuple采用$\gamma=1$。CNN和普通n-tuple的具体实现使用包含棋盘总和与达标项的效用或残差价值，不能直接把三者的更新目标视为完全相同。

\begin{table}[htbp]
\centering
\caption{三种TD价值表示}\label{tab:td}
\begin{tabularx}{\linewidth}{@{}p{0.19\linewidth}YY@{}}\toprule
方法 & 价值表示 & 主要特点\\\midrule
乐观5-tuple & 3张$16^5$权重表，24个对称特征 & 初始总价值设为300，以偏高估值鼓励探索\cite{optimistic}\\\addlinespace[3bp]
CNN-TD & 11个one-hot输入平面，79,874个参数 & 学习棋盘效用残差，使用Adam与梯度裁剪\\\addlinespace[3bp]
n-tuple TD & 行、列、$2\times2$局部模式 & 对称规范化，结合启发式先验与表项更新\\\bottomrule
\end{tabularx}
\end{table}

乐观初始化的作用在于：较少访问的局面仍保留较高估计，智能体有动力尝试这些局面；随着训练，估计逐渐接近实际回报。局部模式表则通过共享相似棋盘的特征，提高经验利用率\cite{optimistic,ntuple}。

\subsection{PPO策略学习}
PPO直接学习动作概率，同时训练价值网络，通过GAE估计优势，并用剪切概率比限制一次更新的幅度\cite{ppo}。本实验采用卷积残差网络，在输出端屏蔽非法动作。评测时冻结权重并选择概率最大的合法动作。与查表TD相比，PPO依赖神经网络训练，其效果还受样本量、奖励尺度和折扣因子影响。

\section{实验设置与超参数}
各方法使用种子10000至10019进行20局测试，运行到自然终局，不设置测试步数上限。测试期间不更新权重，搜索随机数与环境落块随机数分开。相同种子能够统一随机数来源，但策略改变后，棋盘空格不同，实际落块位置也会不同。

表\ref{tab:params}列出实际使用的预算与主要参数。训练局数、环境步数、搜索模拟次数分别对应不同开销，不能简单等同。学习方法的参数以模型保存的记录为依据。

\begin{table}[htbp]
\centering
\caption{计算预算与主要超参数}\label{tab:params}
\begin{tabularx}{\linewidth}{@{}p{0.19\linewidth}p{0.26\linewidth}Y@{}}\toprule
方法 & 预算 & 主要参数\\\midrule
MCTS & 每步128次模拟 & 深度12，$c=70$，随机模拟比例0.15，空格权重18\\\addlinespace[3bp]
完整Expectimax & 搜索至终局并缓存 & 无深度截断，终局达标奖励256\\\addlinespace[3bp]
乐观5-tuple TD & 300,000局 & 训练种子42，初始总价值300，24个特征\\\addlinespace[3bp]
CNN-TD & 50,000局 & 学习率$3\times10^{-4}$，批量1024，探索率0.10降至0.01\\\addlinespace[3bp]
n-tuple TD & 5,000局 & $\alpha=0.2$，探索率0.20降至0.02，空格权重18\\\addlinespace[3bp]
PPO & 50,176个RL环境步 & 学习率$10^{-4}$，$\gamma=0.999$，GAE $\lambda=0.95$，clip=0.2\\\bottomrule
\end{tabularx}
\end{table}

CNN和PPO已有模型的学习率分别为$3\times10^{-4}$和$10^{-4}$，PPO折扣因子为0.999，与当前训练入口的默认值存在差别。PPO还存在教师预训练记录，因此50,176步仅表示强化学习阶段的计数。最终榜单中不同TD模型的训练量、表示方式和先验同时变化，不作单因素消融。下面使用另有记录的历史对照讨论参数影响。

\subsection{超参数对照与讨论}
为讨论参数影响，另整理早期探索记录。这部分与后文20局最终榜单分开：表\ref{tab:hyper}各模型均从头训练10万局，训练种子为42，测试种子为20000至20499，共500局；前70\%训练使用$\alpha=0.1$，之后为0.03，前80\%使用$\epsilon=0.02$，之后为0.002。训练局数和参数日程一致，但不同策略的交互步数与耗时仍有差异。

\begin{table}[htbp]
\centering
\caption{历史TD对照：相同训练局数下的参数与结构变化}\label{tab:hyper}
\begin{tabularx}{\linewidth}{@{}Yrrp{0.24\linewidth}@{}}\toprule
设置 & 平均总和 & 达到256 & 相对基准的变化\\\midrule
对称5-tuple，初值0（基准） & 605.76 & 89.0\% & ---\\
5-tuple，不共享对称 & 417.69 & 72.2\% & 对称特征\\
对称4-tuple & 432.88 & 74.0\% & 局部模式结构\\
对称5-tuple，初值300 & 540.78 & 94.4\% & 乐观初始化\\
对称5-tuple，合并分奖励 & 592.64 & 88.4\% & 奖励目标及尺度\\
对称5-tuple，首次256加奖 & 482.92 & 86.8\% & 奖励增加达标项\\
对称6-tuple，首次256加奖 & 401.52 & 66.6\% & 对比上一行的模式\\\bottomrule
\end{tabularx}
\end{table}

\textbf{表示与初始化。}对称共享使平均总和由417.69提高到605.76；它使旋转、镜像局面共享经验，是本轮较明确的改进。4-tuple低于5-tuple，而采用达标奖励的6-tuple也低于同奖励的5-tuple。由于模式数量和表容量同时变化，这说明当前训练量下更大的局部模式未必更好，不能单独归因于特征长度。初值300将达标率从89.0\%提高到94.4\%，平均总和却降低约65，说明更容易达到256与继续累积总和并不总是一致；现有记录没有比较多个非零初值，不能认定300最优。

\textbf{训练量、学习率与探索率。}表\ref{tab:budget}比较10万局与另一次从头训练30万局的结果，测试仍使用同一组500个种子。三组平均总和均上升，但乐观版本的达标率下降，更长训练没有让两项指标同步提高。两轮的退火时间点也随总预算后移，不能视为同一训练轨迹上的两个检查点。

\begin{table}[!htb]
\centering
\caption{训练预算变化：平均总和与达标率}\label{tab:budget}
\begin{tabularx}{\linewidth}{@{}Yrrrr@{}}\toprule
设置 & 10万局总和 & 30万局总和 & 10万局达标 & 30万局达标\\\midrule
普通对称5-tuple & 605.76 & 629.84 & 89.0\% & 93.0\%\\
乐观对称5-tuple & 540.78 & 618.68 & 94.4\% & 90.4\%\\
5-tuple加达标奖励 & 482.92 & 595.00 & 86.8\% & 90.6\%\\\bottomrule
\end{tabularx}
\end{table}

训练日志中降低学习率与探索率后回报上升，但日志包含探索动作，又与训练进度重合，无法分离各因素的贡献。在另一组从同一权重续训20万局的对照中，固定$\alpha=0.01$，$\epsilon=0.0005$与0分别取得708.79和699.23的平均总和；300局配对bootstrap的差值95\%区间为$[-16.4,35.3]$，尚不足以确认保留微量探索的优势。

\textbf{前瞻深度与计算成本。}固定10万局普通5-tuple权重，只改变决策时的前瞻深度，结果见表\ref{tab:depth}。深度3相对深度2提高约6.7\%的总和，耗时约为14.4倍。该结果支持浅层搜索与TD估值结合；计时来自历史编译实现，排除了编译预热，不能代表当前Python入口速度或完整Expectimax的成本。

\begin{table}[!htb]
\centering
\caption{固定TD权重的搜索深度对照（500局）}\label{tab:depth}
\begin{tabularx}{\linewidth}{@{}Yrrr@{}}\toprule
决策方式 & 平均总和 & 达到256 & 每步耗时/ms\\\midrule
直接TD估值（深度1） & 605.76 & 89.0\% & 0.00076\\
TD与深度2搜索 & 670.88 & 94.2\% & 0.00615\\
TD与深度3搜索 & 716.07 & 95.8\% & 0.08855\\\bottomrule
\end{tabularx}
\end{table}

\textbf{结论边界。}现有证据支持对称共享，以及在允许更多决策时间时加深前瞻。乐观初始化和达标奖励需要同时观察总和与达标率；合并分奖励还改变了奖励尺度，不能只解释为目标差异。学习率、MCTS的模拟次数与探索系数、PPO的clip和折扣因子都缺少逐项多取值对照，目前只能说明所用配置及其作用。以上每种配置只有一个训练种子，评测置信区间也不包含重新训练的波动，不能称为全局最优超参数。


\section{结果与分析}
表\ref{tab:results}给出平均终局总和、生成256的比例和平均步数。MCTS、完整Expectimax、乐观TD及PPO的数据来自逐局记录，CNN与普通n-tuple的数据来自同一20局协议的汇总结果。

\begin{table}[!htb]
\centering
\caption{20局测试结果}\label{tab:results}
\begin{tabularx}{\linewidth}{@{}p{0.30\linewidth}ZZZ@{}}\toprule
方法 & 平均总和 & 达到256 & 平均步数\\\midrule
完整Expectimax & 761.00 & 20/20（100\%） & 344.55\\
乐观5-tuple TD & 612.10 & 20/20（100\%） & 277.60\\
CNN-TD & 418.80 & 16/20（80\%） & 188.50\\
n-tuple TD & 384.00 & 14/20（70\%） & 172.15\\
MCTS & 266.00 & 6/20（30\%） & 119.45\\
PPO & 211.40 & 2/20（10\%） & 94.55\\\bottomrule
\end{tabularx}
\end{table}

\begin{figure}[!htb]
\centering% Native PGFPlots figure: no embedded table or chart screenshot.
\begin{tikzpicture}
\begin{groupplot}[
 group style={group size=2 by 1,horizontal sep=12mm},
 width=0.45\linewidth,height=59mm,scale only axis=false,
 tick label style={font=\fontsize{9bp}{10bp}\selectfont},
 title style={font=\fontsize{10bp}{12bp}\selectfont},
 label style={font=\fontsize{10bp}{12bp}\selectfont},
 axis lines=left,ymajorgrids,grid style={gray!20},
 xticklabel style={rotate=35,anchor=east},
]
\nextgroupplot[title={A. 平均终局总和},ymin=0,ymax=900,
 ylabel={棋盘总和},xtick={1,2,3,4,5,6},
 xticklabels={MCTS,Expectimax,TD-Opt5,TD-CNN,TD-nTuple,PPO},
 xmin=0.4,xmax=6.6,ybar,bar width=13pt,
 nodes near coords,nodes near coords style={font=\fontsize{8bp}{9bp}\selectfont},
 every node near coord/.append style={/pgf/number format/fixed,/pgf/number format/precision=1}]

\addplot+[draw=none,bar shift=0pt,fill=gray!65,nodes near coords style={text=black}] coordinates {(1,266)};
\addplot+[draw=none,bar shift=0pt,fill=violet!65,nodes near coords style={text=black}] coordinates {(2,761)};
\addplot+[draw=none,bar shift=0pt,fill=orange!70!brown,nodes near coords style={text=black}] coordinates {(3,612.1)};
\addplot+[draw=none,bar shift=0pt,fill=teal!65,nodes near coords style={text=black}] coordinates {(4,418.8)};
\addplot+[draw=none,bar shift=0pt,fill=violet!40,nodes near coords style={text=black}] coordinates {(5,384)};
\addplot+[draw=none,bar shift=0pt,fill=red!45!gray,nodes near coords style={text=black}] coordinates {(6,211.4)};
\nextgroupplot[title={B. 逐局总和分布},ymin=0,ymax=1100,
 xmin=0.4,xmax=4.6,xtick={1,2,3,4},xticklabels={MCTS,Expectimax,TD-Opt5,PPO},
 boxplot/draw direction=y,boxplot/box extend=0.5]

\addplot+[fill=gray!65,draw=black,solid,mark=o,boxplot prepared={draw position=1,lower whisker=120.0,lower quartile=226.0,median=259.0,upper quartile=309.5,upper whisker=394.0}] coordinates {(0,466.0)};
\addplot[only marks,mark=diamond*,mark options={fill=white,draw=black},mark size=2pt] coordinates {(1,266.0)};
\addplot+[fill=violet!65,draw=black,solid,mark=o,boxplot prepared={draw position=2,lower whisker=500.0,lower quartile=513.5,median=833.0,upper quartile=995.0,upper whisker=1016.0}] coordinates {};
\addplot[only marks,mark=diamond*,mark options={fill=white,draw=black},mark size=2pt] coordinates {(2,761.0)};
\addplot+[fill=orange!70!brown,draw=black,solid,mark=o,boxplot prepared={draw position=3,lower whisker=392.0,lower quartile=509.0,median=514.0,upper quartile=753.0,upper whisker=948.0}] coordinates {};
\addplot[only marks,mark=diamond*,mark options={fill=white,draw=black},mark size=2pt] coordinates {(3,612.1)};
\addplot+[fill=red!45!gray,draw=black,solid,mark=o,boxplot prepared={draw position=4,lower whisker=46.0,lower quartile=144.0,median=226.0,upper quartile=262.5,upper whisker=368.0}] coordinates {(0,508.0)};
\addplot[only marks,mark=diamond*,mark options={fill=white,draw=black},mark size=2pt] coordinates {(4,211.4)};
\end{groupplot}\end{tikzpicture}

\caption{平均终局总和与逐局分布。右图仅包含有逐局记录的四种方法，箱体表示四分位范围，横线为中位数，菱形为均值。}\label{fig:results}
\end{figure}

完整Expectimax的平均总和最高，比乐观TD高148.9，约24.3\%。两者均在20局中生成256，完整搜索在较高总和局面上的收益更明显。乐观TD排名第二，推理主要是特征查表，适合在线计算资源有限的场景。CNN与普通n-tuple都超过当前MCTS，但由于训练量与先验不同，尚不能据此判断某种表示本身更优。

MCTS的结果说明，小预算随机模拟配合启发式估值，在3×3的紧凑棋盘上仍容易陷入死局。PPO的平均总和最低，当前训练量与奖励设置下尚未形成稳定策略。其原因需要进一步训练和对照实验确认，不能仅凭这组结果归因于PPO算法本身。平均步数与总和总体一致，这是因为活得更久通常能够接纳更多出生块。

20局全达标仍是有限样本结果。若把各局视为独立同分布试验，20/20对应的95\% Wilson区间约为[83.9\%，100\%]。在这20个种子上进行20,000次配对bootstrap，Expectimax与乐观TD平均总和差的95\%区间约为[23.6，273.0]，该区间不包含重新训练带来的波动。另一次乐观TD的500局测试平均总和为639.408，463局生成256；它使用不同测试设置，不与表\ref{tab:results}合并。

\section{总结}
在本实验的3×3环境中，完整Expectimax取得最高平均总和，乐观5-tuple TD则兼顾效果与在线决策开销。搜索能充分利用已知转移规则，TD可以把搜索或交互经验转化为快速估值，二者各有适用条件。后续可在共同奖励目标下统一运行时间或每步时限，分别记录训练、建表、内存与推理开销，并用多个训练种子和新的测试种子重复实验。对学习率、探索率和搜索预算逐项调整，才能更清楚地区分参数设置与算法结构的影响。

\section{小组分工}
\noindent\begin{tabularx}{\linewidth}{@{}p{0.13\linewidth}Y@{}}\toprule
成员 & 负责内容\\\midrule
蔡明辰 & 环境适配、随机种子与规则测试，统筹代码整合\\\addlinespace[3bp]
潘亮 & TD方法实现、模型训练与超参数整理\\\addlinespace[3bp]
陈达 & MCTS、Expectimax和PPO的方法实现及对照分析\\\addlinespace[3bp]
孟维祎 & 评测数据整理、图表制作与报告撰写\\\bottomrule
\end{tabularx}

\begin{thebibliography}{9}
\setlength{\itemsep}{2bp}\setlength{\parskip}{0bp}
\bibitem{optimistic} Guei H, Chen L-P, Wu I-C. \href{https://arxiv.org/abs/2111.11090}{Optimistic Temporal Difference Learning for 2048}. 2021.
\bibitem{ntuple} Szubert M, Jaskowski W. \href{https://mszubert.github.io/papers/Szubert_2014_CIG.pdf}{Temporal Difference Learning of N-Tuple Networks for the Game 2048}. CIG, 2014.
\bibitem{ppo} Schulman J, et al. \href{https://arxiv.org/abs/1707.06347}{Proximal Policy Optimization Algorithms}. 2017.
\end{thebibliography}
\end{document}
